\chapter{无界自伴算子的谱定理}
\label{chap:unbounded-spectral-theorem}

本章给出并证明无界自伴算子的谱定理.投影值测度、直积分和乘法算子三种形式与有界情形基本一致,差别主要在于必须描述算子的定义域.证明的核心是把无界自伴算子 \(A\) 化为有界酉算子
\[
U=(A+iI)(A-iI)^{-1},
\]
再对 \(U\) 使用有界正规算子的谱定理.为此,先证明有界正规算子的谱定理.第~\ref{sec:one-parameter-unitary-groups} 节还讨论强连续一参数酉群及 Stone 定理.

\section{谱定理的陈述}
\label{sec:unbounded-spectral-statements}

设 \(\mu\) 是 \((X,\Omega)\) 上取值于 \(\mathcal{B}(\mathcal{H})\) 的投影值测度.对 \(\psi\in\mathcal{H}\),定义正实测度 \(\mu_{\psi}(E)=\langle\psi,\mu(E)\psi\rangle\).对于有界可测函数 \(f\),由积分的乘法性和共轭对应于伴随,有
\[
\left\|\left(\int_{X}f\,\mathrm{d}\mu\right)\psi\right\|^{2}
=\int_{X}|f|^{2}\,\mathrm{d}\mu_{\psi}.
\label{eq:unbounded-pvm-norm-identity}
\]
若 \(f\) 无界,仍可把右端有限的 \(\psi\) 作为积分算子的定义域.

\begin{proposition}[无界投影值积分]\label{prop:unbounded-pvm-integral}
设 \(\mu\) 是 \((X,\Omega)\) 上取值于 \(\mathcal{B}(\mathcal{H})\) 的投影值测度,\(f:X\to\mathbb{C}\) 可测但不必有界.令
\[
W_{f}:=\left\{\psi\in\mathcal{H}:\int_{X}|f(\lambda)|^{2}\,\mathrm{d}\mu_{\psi}(\lambda)<\infty\right\}.
\label{eq:unbounded-pvm-domain}
\]
则存在唯一的定义域为 \(W_{f}\) 的无界算子,记为 \(\int_{X}f\,\mathrm{d}\mu\),使得
\[
\left\langle\psi,\left(\int_{X}f\,\mathrm{d}\mu\right)\psi\right\rangle
=\int_{X}f(\lambda)\,\mathrm{d}\mu_{\psi}(\lambda),
\]
且式~\eqref{eq:unbounded-pvm-norm-identity} 对所有 \(\psi\in W_{f}\) 成立.
\end{proposition}

当 \(f\) 有界时,\(W_{f}=\mathcal{H}\),本命题退化为第~7 章的有界算子值积分.

\begin{proposition}[无界积分的二次型性质]\label{prop:unbounded-integral-quadratic}
令 \(W_{f}\) 如式~\eqref{eq:unbounded-pvm-domain}.则:
\begin{enumerate}
\item \(W_{f}\) 是 \(\mathcal{H}\) 的稠密子空间,且 \(Q_{f}(\psi)=\int f\,\mathrm{d}\mu_{\psi}\) 是定义在 \(W_{f}\) 上的二次型;
\item 关联的半双线性型 \(L_{f}\) 满足
\[
|L_{f}(\phi,\psi)|\leq\|\phi\|\left(\int_{X}|f|^{2}\,\mathrm{d}\mu_{\psi}\right)^{1/2};
\label{eq:unbounded-integral-form-bound}
\]
\item 对每个 \(\psi\in W_{f}\),存在唯一 \(\chi\in\mathcal{H}\),使 \(L_{f}(\phi,\psi)=\langle\phi,\chi\rangle\) 对所有 \(\phi\in W_{f}\) 成立,且
\[
\|\chi\|^{2}=\int_{X}|f|^{2}\,\mathrm{d}\mu_{\psi}.
\]
\end{enumerate}
\end{proposition}

\begin{proof}
\(W_{f}\) 对数乘封闭.若 \(\phi,\psi\in W_{f}\),则利用 \(\mu(E)^{2}=\mu(E)=\mu(E)^{*}\),
\[
\mu_{\phi+\psi}(E)=\|\mu(E)(\phi+\psi)\|^{2}
\leq2\mu_{\phi}(E)+2\mu_{\psi}(E),
\]
所以 \(W_{f}\) 对加法封闭.

令 \(E_{n}=\{x:|f(x)|<n\}\).若 \(\psi\in\operatorname{Range}(\mu(E_{n}))\),则 \(\mu_{\psi}(E_{n}^{c})=0\),从而
\[
\int_{X}|f|^{2}\,\mathrm{d}\mu_{\psi}
=\int_{E_{n}}|f|^{2}\,\mathrm{d}\mu_{\psi}
\leq n^{2}\mu_{\psi}(E_{n})<\infty.
\]
故这些值域都包含于 \(W_{f}\).由于 \(\bigcup E_{n}=X\),这些值域的并稠密,故 \(W_{f}\) 稠密.

若 \(f\) 有界,所有结论由有界积分直接得到;若无界,令 \(f_{n}=f\mathbf{1}_{E_{n}}\),用单调收敛定理得到相应结论.式~\eqref{eq:unbounded-integral-form-bound} 说明固定 \(\psi\) 后 \(\phi\mapsto L_{f}(\phi,\psi)\) 是有界泛函,Riesz 定理给出唯一 \(\chi\).令 \(\phi=\chi\) 并比较两侧,即得范数等式,线性性则来自 \(L_{f}\) 关于 \(\psi\) 的线性.
\end{proof}

\begin{proposition}[实值无界积分的自伴性]\label{prop:real-unbounded-integral-selfadjoint}
若 \(f:X\to\mathbb{R}\) 可测,则 \(\int f\,\mathrm{d}\mu\) 在 \(W_{f}\) 上自伴.
\end{proposition}

\begin{proof}
令 \(F_{n}=\{x:n-1\leq|f(x)|<n\}\),令 \(W^{n}=\operatorname{Range}(\mu(F_{n}))\).每个 \(W^{n}\subset W_{f}\),且积分二次型在 \(W^{n}\) 上有界.若 \(\phi\perp W^{n}\)、\(\psi\in W^{n}\),则 \(\mu_{\phi+\psi}=\mu_{\phi}+\mu_{\psi}\),所以由极化恒等式,\(\int f\,\mathrm{d}\mu\) 把 \(W^{n}\) 映到自身.

在 \(W^{n}\) 上,它与有界函数 \(f\mathbf{1}_{F_{n}}\) 的积分一致,因而是有界自伴算子.Hilbert 空间是这些 \(W^{n}\) 的正交直和;若写 \(\psi=(\psi_{1},\psi_{2},\ldots)\),则
\[
W_{f}=\left\{\psi:\sum_{n}\|A_{n}\psi_{n}\|^{2}<\infty\right\},
\]
其中 \(A_{n}\) 是限制在 \(W^{n}\) 上的有界算子.于是积分算子由 \((A_{1}\psi_{1},A_{2}\psi_{2},\ldots)\) 给出,例~\ref{ex:direct-sum-essential-selfadjoint} 说明它自伴.
\end{proof}

\begin{theorem}[谱定理:无界情形第一种形式]\label{thm:unbounded-spectral-first}
设 \(A\) 是 \(\mathcal{H}\) 上的自伴算子.则存在唯一的、定义在 \(\sigma(A)\subset\mathbb{R}\) 的 Borel \(\sigma\)-代数上的投影值测度 \(\mu_{A}\),使得
\[
A=\int_{\sigma(A)}\lambda\,\mathrm{d}\mu_{A}(\lambda).
\label{eq:unbounded-spectral-first}
\]
这里等式包括定义域相等:左端定义域等于右端积分中函数 \(f(\lambda)=\lambda\) 的定义域 \(W_{f}\).
\end{theorem}

\begin{definition}[无界函数演算]\label{def:unbounded-functional-calculus}
若 \(A\) 自伴,且 \(f\) 是 \(\sigma(A)\) 上任意可测函数,则定义可能无界的算子
\[
f(A):=\int_{\sigma(A)}f(\lambda)\,\mathrm{d}\mu_{A}(\lambda).
\]
\end{definition}

\begin{definition}[无界情形的谱子空间]\label{def:unbounded-spectral-subspace}
对自伴算子 \(A\) 和任意 Borel 集 \(E\subset\mathbb{R}\),定义
\[
V_{E}:=\operatorname{Range}(\mu_{A}(E)).
\]
\end{definition}

\begin{definition}[测量概率]\label{def:unbounded-measurement-probability}
若 \(\psi\) 是单位向量,则定义 \(A\) 在状态 \(\psi\) 下的概率测度
\[
\mu_{\psi}^{A}(E):=\langle\psi,\mu_{A}(E)\psi\rangle.
\]
它描述测量可观测量 \(A\) 时所得结果的概率分布.
\end{definition}

\begin{proposition}[无界谱子空间的性质]\label{prop:unbounded-spectral-subspaces}
对自伴算子 \(A\) 的谱子空间 \(V_{E}\):
\begin{enumerate}
\item 若 \(E\) 是有界集,则 \(V_{E}\subset\operatorname{Dom}(A)\),且 \(V_{E}\) 在 \(A\) 下不变,\(A|_{V_{E}}\) 有界;
\item 若 \(E\subset(\lambda_{0}-\varepsilon,\lambda_{0}+\varepsilon)\),则
\[
\|(A-\lambda_{0}I)\psi\|\leq\varepsilon\|\psi\|,qquad \psi\in V_{E}.
\]
\end{enumerate}
\end{proposition}

\begin{proof}
第一点是因为 \(f(\lambda)=\lambda\) 在有界集 \(E\) 上有界,故 \(f(A)\) 在 \(V_{E}\) 上有界.第二点把 \(A|_{V_{E}}\) 与有界函数 \(\lambda\mathbf{1}_{E}(\lambda)\) 的函数演算比较即可.
\end{proof}

\begin{theorem}[谱定理:无界情形第二种形式]\label{thm:unbounded-spectral-direct-integral}
设 \(A\) 是自伴算子.则存在 \(\sigma\)-有限测度 \(\mu\)、直积分
\[
\int_{\sigma(A)}^{\oplus}\mathcal{H}_{\lambda}\,\mathrm{d}\mu(\lambda)
\]
以及从 \(\mathcal{H}\) 到该直积分的酉映射 \(U\),使
\[
U(\operatorname{Dom}(A))
=\left\{s:\int_{\sigma(A)}\|\lambda s(\lambda)\|_{\lambda}^{2}\,\mathrm{d}\mu(\lambda)<\infty\right\}
\]
且
\[
\bigl[UAU^{-1}s\bigr](\lambda)=\lambda s(\lambda)
\]
对所有属于该定义域的截面成立.
\end{theorem}

\begin{theorem}[谱定理:无界乘法算子形式]\label{thm:unbounded-spectral-multiplication}
设 \(A\) 自伴.则存在 \(\sigma\)-有限测度空间 \((X,\mu)\)、可测实值函数 \(h\) 和酉映射 \(U:\mathcal{H}\to L^{2}(X,\mu)\),使
\[
U(\operatorname{Dom}(A))=\{\psi\in L^{2}(X,\mu):h\psi\in L^{2}(X,\mu)\}
\]
且
\[
\bigl[UAU^{-1}\psi\bigr](x)=h(x)\psi(x).
\]
\end{theorem}

两定理的证明见第~\ref{sec:proof-unbounded-spectral} 节.

\section{Stone 定理与一参数酉群}
\label{sec:one-parameter-unitary-groups}

\begin{definition}[一参数酉群]\label{def:one-parameter-unitary-group}
Hilbert 空间上的一参数酉群是酉算子族 \(U(t)\)(\(t\in\mathbb{R}\)),满足
\[
U(0)=I,\qquad U(s+t)=U(s)U(t).
\]
若对所有 \(\psi\in\mathcal{H}\) 及 \(t\in\mathbb{R}\),
\[
\lim_{s\to t}\|U(s)\psi-U(t)\psi\|=0,
\label{eq:strong-continuity-unitary-group}
\]
则称其强连续.
\end{definition}

\begin{example}[平移一参数酉群]\label{ex:translation-unitary-group}
在 \(L^{2}(\mathbb{R}^{n})\) 上令
\[
(U_{a}(t)\psi)(x)=\psi(x+ta).
\label{eq:translation-group}
\]
则 \(U_{a}(t)\) 是强连续一参数酉群.
\end{example}

\begin{proof}
酉性和群性质直接可见.若 \(\psi\) 连续且紧支集,则在紧集上函数一致连续,所以 \(\psi(x+ta)\) 一致趋于 \(\psi(x)\),从而在 \(L^{2}\) 中收敛.再利用 \(C_{c}(\mathbb{R}^{n})\) 在 \(L^{2}\) 中稠密及酉性,即可推广到所有 \(\psi\).
\end{proof}

若 \(a\neq0\),这个群不具有算子范数连续性.取支集在足够小球中的非零函数 \(\psi\),则对任意 \(t\neq0\) 可使 \(U_{a}(t)\psi\perp\psi\),从而
\[
\|U_{a}(t)-I\|\geq\sqrt{2}.
\]

\begin{definition}[无穷小生成元]\label{def:unitary-generator}
强连续一参数酉群 \(U(t)\) 的无穷小生成元是算子 \(A\):
\[
A\psi=\lim_{t\to0}\frac{1}{i}\frac{U(t)\psi-\psi}{t},
\label{eq:unitary-generator}
\]
其定义域由使该极限在 \(\mathcal{H}\) 的范数拓扑中存在的 \(\psi\) 组成.
\end{definition}

\begin{proposition}[自伴算子生成一参数酉群]\label{prop:selfadjoint-generates-unitary}
设 \(A\) 自伴,令 \(U(t)=e^{itA}\),其中指数由函数演算定义.则:
\begin{enumerate}
\item \(U(t)\) 是强连续一参数酉群;
\item 对所有 \(\psi\in\operatorname{Dom}(A)\),
\[
A\psi=\lim_{t\to0}\frac{1}{i}\frac{U(t)\psi-\psi}{t};
\]
\item 若对某个 \(\psi\in\mathcal{H}\) 上述极限存在,则 \(\psi\in\operatorname{Dom}(A)\),且极限为 \(A\psi\).
\end{enumerate}
\end{proposition}

\begin{proof}
由于 \(\sigma(A)\subset\mathbb{R}\),函数 \(e^{it\lambda}\) 有界且模为 1,所以函数演算给出酉算子;乘法性给出群性质.对强连续性,令 \(\mu_{\psi}^{A}\) 为谱测度,则
\[
\|U(t)\psi-U(s)\psi\|^{2}
=\int_{\mathbb{R}}|e^{it\lambda}-e^{is\lambda}|^{2}\,\mathrm{d}\mu_{\psi}^{A}(\lambda),
\]
由支配收敛趋于零.

若 \(\psi\in\operatorname{Dom}(A)\),则
\[
\left\|\frac{1}{i}\frac{U(t)\psi-\psi}{t}-A\psi\right\|^{2}
=\int\left|\frac{e^{it\lambda}-1}{it}-\lambda\right|^{2}\,\mathrm{d}\mu_{\psi}^{A}(\lambda),
\]
被 \(4\lambda^{2}\) 控制,故由支配收敛趋于零.第三点可令极限存在的向量组成的算子记为 \(B\),证明 \(B\) 对称;第二点说明 \(B\) 延拓 \(A\),而对称算子无严格对称延拓(练习~\ref{ex:extension-adjoint-properties}),故 \(B=A\).
\end{proof}

\begin{theorem}[Stone 定理]\label{thm:stone}
若 \(U(t)\) 是强连续一参数酉群,则其无穷小生成元 \(A\) 稠密定义且自伴,并且
\[
U(t)=e^{itA}qquad(t\in\mathbb{R}).
\]
\end{theorem}

\begin{lemma}[生成元定义域的不变性]\label{lem:generator-domain-invariant}
若 \(\psi\in\operatorname{Dom}(A)\),则对所有 \(t\),\(U(t)\psi\in\operatorname{Dom}(A)\),且
\[
\lim_{h\to0}\frac{U(t+h)\psi-U(t)\psi}{h}
=iU(t)A\psi=iAU(t)\psi.
\]
\end{lemma}

\begin{proof}
由群性质,左端等于 \(U(t)[U(h)\psi-\psi]/h\),极限为 \(iU(t)A\psi\);另一方面把它写成 \([U(h)U(t)\psi-U(t)\psi]/h\),由生成元定义得到另一种表达式.
\end{proof}

\begin{lemma}[生成元的稠密定义域]\label{lem:generator-dense}
强连续一参数酉群的生成元稠密定义.
\end{lemma}

\begin{proof}
对紧支集连续函数 \(f\),定义有界算子
\[
B_{f}:=\int_{-\infty}^{\infty}f(\tau)U(\tau)\,\mathrm{d}\tau,
\]
其弱意义由
\[
\langle\phi,B_{f}\psi\rangle
=\int f(\tau)\langle\phi,U(\tau)\psi\rangle\,\mathrm{d}\tau
\]
给出.群性质和变量替换表明 \(B_{f}\psi\in\operatorname{Dom}(A)\).取非负函数 \(f_{n}\in C_{c}^{\infty}\),支集在 \([-1/n,1/n]\) 且积分为 1,则
\[
\|B_{f_{n}}\psi-\psi\|
\leq\sup_{|\tau|\leq1/n}\|U(\tau)\psi-\psi\|\to0.
\]
所以生成元定义域稠密.
\end{proof}

\begin{proof}[Stone 定理的证明]
由引理~\ref{lem:generator-dense},生成元 \(A\) 稠密定义;如命题~\ref{prop:selfadjoint-generates-unitary} 的证明,\(A\) 对称.

若 \(\psi\in\ker(A^{*}-iI)\),取 \(\phi\in\operatorname{Dom}(A)\),令 \(y(t)=\langle U(t)\phi,\psi\rangle\).由引理~\ref{lem:generator-domain-invariant},
\[
y'(t)=\langle iAU(t)\phi,\psi\rangle
=\langle iU(t)\phi,A^{*}\psi\rangle=y(t).
\]
所以 \(y(t)=y(0)e^{t}\).另一方面 \(|y(t)|\leq\|\phi\|\|\psi\|\),故 \(y(0)=0\),即 \(\psi\perp\operatorname{Dom}(A)\),从而 \(\psi=0\).同理 \(\ker(A^{*}+iI)=\{0\}\),由推论~\ref{cor:essential-selfadjoint-criterion},\(A\) 本质自伴.

令 \(V(t)=e^{itA^{\mathrm{cl}}}\).对 \(\psi\in\operatorname{Dom}(A)\),令 \(w(t)=U(t)\psi-V(t)\psi\).由引理~\ref{lem:generator-domain-invariant},\(w'(t)=iAw(t)\),且 \(w(0)=0\).因为 \(A\) 对称,\(\mathrm{d}\|w(t)\|^{2}/\mathrm{d}t=0\),故 \(w(t)=0\).稠密性和酉性推出 \(U(t)=V(t)\) 对所有向量成立.再由命题~\ref{prop:selfadjoint-generates-unitary} 的第 2、3 点,原生成元等于 \(A^{\mathrm{cl}}\),所以 \(A\) 自伴,且 \(U(t)=e^{itA}\).
\end{proof}

\section{有界正规算子的谱定理}
\label{sec:bounded-normal-spectral-theorem}

\begin{definition}[正规算子]\label{def:normal-operator}
若有界算子 \(A\) 与其伴随相交换,即 \(AA^{*}=A^{*}A\),则称 \(A\) 为正规算子.
\end{definition}

自伴算子、斜自伴算子(\(A^{*}=-A\))和酉算子都是正规算子.正规算子的谱不必包含于实轴,但仍是复平面中的非空、闭、有界集合.

\begin{theorem}[正规算子的谱定理]\label{thm:bounded-normal-spectral}
若 \(A\in\mathcal{B}(\mathcal{H})\) 正规,则存在唯一的投影值测度 \(\mu_{A}\)(定义在 \(\sigma(A)\) 的 Borel \(\sigma\)-代数上),使
\[
\int_{\sigma(A)}\lambda\,\mathrm{d}\mu_{A}(\lambda)=A.
\]
此外,对任意可测集 \(E\),\(\operatorname{Range}(\mu_{A}(E))\) 在 \(A\) 和 \(A^{*}\) 下都不变.
\end{theorem}

其直积分形式和乘法算子形式与自伴情形完全相同,只是不再要求谱位于实轴,也不再要求乘法函数为实值.

\begin{proposition}[正规算子的范数等于谱半径]\label{prop:normal-norm-spectral-radius}
若 \(A\) 正规,则 \(\|A\|=R(A)\).
\end{proposition}

\begin{lemma}[可交换乘积的谱半径]\label{lem:commuting-spectral-radius-product}
若 \(A,B\in\mathcal{B}(\mathcal{H})\) 相交换,则
\[
R(AB)\leq R(A)R(B).
\]
\end{lemma}

\begin{proof}
对任意 \(T>R(A)\),引理~\ref{lem:selfadjoint-spectral-radius} 证明中的预解级数论证给出 \(\|A^{m}\|/T^{m}\to0\).若 \(S>R(A)\)、\(T>R(B)\),则
\[
\frac{\|(AB)^{m}\|}{S^{m}T^{m}}
\leq\frac{\|A^{m}\|}{S^{m}}\frac{\|B^{m}\|}{T^{m}}\to0.
\]
又因 \(A,B\) 相交换,预解展开
\[
(AB-\lambda I)^{-1}=-\sum_{m=0}^{\infty}\frac{A^{m}B^{m}}{\lambda^{m+1}}
\]
在 \(|\lambda|>R(A)R(B)\) 时收敛.因此该区域属于 \(AB\) 的预解集.
\end{proof}

\begin{proof}[命题~\ref{prop:normal-norm-spectral-radius} 的证明]
一般总有 \(\|A\|\geq R(A)\).另一方面 \(\|A\|^{2}=\|A^{*}A\|\),且 \(A^{*}A\) 自伴.由引理~\ref{lem:commuting-spectral-radius-product},
\[
\|A\|^{2}=R(A^{*}A)\leq R(A^{*})R(A)\leq\|A^{*}\|R(A)=\|A\|R(A).
\]
除非 \(A=0\),约去 \(\|A\|\) 得 \(\|A\|\leq R(A)\).
\end{proof}

\begin{theorem}[正规算子的多项式谱映射]\label{thm:normal-polynomial-spectral-mapping}
若 \(A\) 正规,且 \(p\) 是两个变量的多项式,则
\[
\sigma(p(A,A^{*}))
=\{p(\lambda,\overline\lambda):\lambda\in\sigma(A)\}.
\label{eq:normal-polynomial-spectral-mapping}
\]
\end{theorem}

若 \(p(\lambda,\overline\lambda)=\lambda^{m}\overline\lambda^{n}\),则 \(p(A,A^{*})=A^{m}(A^{*})^{n}\).正规性保证相应代数同态的乘法性.

\begin{definition}[\(\varepsilon\)-近似特征向量]\label{def:epsilon-almost-eigenvector}
若 \(\psi\neq0\) 且
\[
\|(A-\lambda I)\psi\|<\varepsilon\|\psi\|,
\]
则称 \(\psi\) 是特征值 \(\lambda\) 对应的 \(\varepsilon\)-近似特征向量.
\end{definition}

\begin{lemma}[正规算子的近似特征向量]\label{lem:normal-almost-eigenvector}
设 \(A\) 正规.
\begin{enumerate}
\item 若 \(\psi\) 是 \(A\) 对应 \(\lambda\) 的 \(\varepsilon\)-近似特征向量,则它也是 \(A^{*}\) 对应 \(\overline\lambda\) 的 \(\varepsilon\)-近似特征向量;
\item \(\lambda\in\sigma(A)\),当且仅当对每个 \(\varepsilon>0\),都存在相应的 \(\varepsilon\)-近似特征向量.
\end{enumerate}
\end{lemma}

\begin{proof}
正规性给出恒等式
\[
\|(A^{*}-\overline\lambda I)\psi\|^{2}
=\|(A-\lambda I)\psi\|^{2}.
\]
这证明第一点.若每个 \(\varepsilon\) 都存在近似特征向量,则 \(A-\lambda I\) 不可能有有界逆.反之,若某个 \(\varepsilon\) 不存在,则 \(\|(A-\lambda I)\psi\|\geq\varepsilon\|\psi\|\),同一估计对 \(A^{*}-\overline\lambda I\) 成立;由值域与伴随核的关系和值域闭性,\(A-\lambda I\) 可逆且逆有界.
\end{proof}

\begin{lemma}[近似特征向量的多项式估计]\label{lem:normal-polynomial-almost}
对正规算子 \(A\)、两个变量的多项式 \(p\) 及 \(\lambda\in\mathbb{C}\),存在常数 \(C\),使得若 \(\psi\) 是 \(A\) 对应 \(\lambda\) 的 \(\varepsilon\)-近似特征向量,则 \(\psi\) 是 \(p(A,A^{*})\) 对应 \(p(\lambda,\overline\lambda)\) 的 \(C\varepsilon\)-近似特征向量.
\end{lemma}

\begin{proof}
把 \(p(A,A^{*})-p(\lambda,\overline\lambda)I\) 分解为 \(A^{k}(A^{*})^{l}-\lambda^{k}\overline\lambda^{l}\) 的线性组合,并对 \(k+l\) 归纳.归纳步骤使用
\[
\bigl[A^{k}(A^{*})^{l}-\lambda^{k}\overline\lambda^{l}I\bigr]\psi
=A^{k-1}(A^{*})^{l}(A-\lambda I)\psi
+\lambda\bigl[A^{k-1}(A^{*})^{l}-\lambda^{k-1}\overline\lambda^{l}I\bigr]\psi.
\]
\end{proof}

\begin{lemma}[正规算子的谱子空间近似]\label{lem:normal-spectral-subspace-almost}
设 \(A\) 正规,\(p\) 是两个变量的多项式,\(\mu\in\sigma(p(A,A^{*}))\).对每个 \(\varepsilon>0\),存在非零闭子空间 \(W_{\varepsilon}\),使其在 \(A,A^{*}\) 下不变,且其中每个非零元素都是 \(p(A,A^{*})\) 对应 \(\mu\) 的 \(\varepsilon\)-近似特征向量.
\end{lemma}

\begin{proof}
令 \(B=p(A,A^{*})-\mu I\).则 \(B\) 正规且 \(0\in\sigma(B)\).由上一个引理,0 属于自伴算子 \(B^{*}B\) 的谱.对 \(B^{*}B\) 使用有界自伴谱定理,令 \(W_{\varepsilon}\) 为区间 \([0,\varepsilon^{2})\) 对应的谱子空间.它非零,在 \(B^{*}B\) 下不变,且
\[
\langle B\psi,B\psi\rangle
=\langle\psi,B^{*}B\psi\rangle\leq\varepsilon^{2}\|\psi\|^{2}.
\]
由于 \(A,A^{*}\) 与 \(B^{*}B\) 相交换,它们保持该谱子空间.
\end{proof}

\begin{proof}[定理~\ref{thm:normal-polynomial-spectral-mapping} 的证明]
若 \(\lambda\in\sigma(A)\),引理~\ref{lem:normal-almost-eigenvector} 给出任意精度的近似特征向量,再由引理~\ref{lem:normal-polynomial-almost},\(p(\lambda,\overline\lambda)\in\sigma(p(A,A^{*}))\).

反过来,若 \(\mu\in\sigma(p(A,A^{*}))\),取引理~\ref{lem:normal-spectral-subspace-almost} 中的 \(W_{\varepsilon}\).限制 \(A\) 到该空间仍正规,其谱非空.取 \(\lambda\in\sigma(A|_{W_{\varepsilon}})\),则存在近似特征向量,同时它既是 \(p(A,A^{*})\) 对应 \(p(\lambda,\overline\lambda)\) 的近似特征向量,又是对应 \(\mu\) 的近似特征向量.因此
\[
|\mu-p(\lambda,\overline\lambda)|<C\varepsilon+\varepsilon.
\]
令 \(\varepsilon\to0\),利用 \(\sigma(A)\) 紧性取收敛子列,得到 \(\mu=p(\lambda,\overline\lambda)\).
\end{proof}

由复 Stone--Weierstrass 定理,\(\lambda,\overline\lambda\) 的多项式在 \(C(\sigma(A);\mathbb{C})\) 中稠密.结合
\[
\|p(A,A^{*})\|=\sup_{\lambda\in\sigma(A)}|p(\lambda,\overline\lambda)|,
\]
可以按与第~8 章相同的步骤构造正规算子的连续函数演算和投影值测度.第二阶段的 Riesz 表示论证完全相同,所以得到定理~\ref{thm:bounded-normal-spectral}.

\section{无界自伴谱定理的证明}
\label{sec:proof-unbounded-spectral}

定义 Cayley 变换
\[
C(x):=\frac{x+i}{x-i},\qquad x\in\mathbb{R}.
\label{eq:cayley-map}
\]
它把实轴一一映到单位圆去掉点 1 的部分,逆映射为
\[
D(u):=i\frac{u+1}{u-1},qquad u\in S^{1}\setminus\{1\}.
\label{eq:cayley-inverse}
\]
且 \(C(x)\to1\) 当 \(x\to\pm\infty\).

\begin{theorem}[Cayley 变换]\label{thm:cayley-transform}
设 \(A\) 自伴,令
\[
U\psi=(A+iI)(A-iI)^{-1}\psi.
\]
则:
\begin{enumerate}
\item \(U\) 是酉算子;
\item \(U-I\) 是单射;
\item \(\operatorname{Range}(U-I)=\operatorname{Dom}(A)\),并且对该值域中的 \(\psi\),
\[
A\psi=i(U+I)(U-I)^{-1}\psi.
\label{eq:cayley-recovery}
\]
\end{enumerate}
\end{theorem}

\begin{proof}
\((A-iI)^{-1}\) 单射且把 \(\mathcal{H}\) 映到 \(\operatorname{Dom}(A)\);\(A+iI\) 单射地把该定义域映到 \(\mathcal{H}\),故 \(U\) 双射.对 \(\phi\in\operatorname{Dom}(A)\),交叉项相消给出
\[
\|(A+iI)\phi\|^{2}=\|A\phi\|^{2}+\|\phi\|^{2}=\|(A-iI)\phi\|^{2}.
\]
令 \(\phi=(A-iI)^{-1}\psi\),得 \(\|U\psi\|=\|\psi\|\),故 \(U\) 酉.

直接计算
\[
U-I=2i(A-iI)^{-1}.
\label{eq:cayley-difference}
\]
由逆算子的单射性,\(U-I\) 单射;由其值域等于 \((A-iI)^{-1}\) 的值域,\(\operatorname{Range}(U-I)=\operatorname{Dom}(A)\).最后由
\[
(U+I)(U-I)^{-1}\psi=\frac{1}{2i}[(A+iI)+(A-iI)]\psi
=\frac{1}{i}A\psi
\]
得到式~\eqref{eq:cayley-recovery}.
\end{proof}

对 \(U\) 使用正规算子谱定理,并通过 \(D\) 把测度从单位圆转移回实轴.

\begin{proposition}[从 Cayley 变换恢复算子]\label{prop:cayley-functional-calculus}
设 \(U\) 是定理~\ref{thm:cayley-transform} 中的酉算子,\(D\) 如式~\eqref{eq:cayley-inverse}.则
\[
A=D(U),
\label{eq:cayley-functional-calculus}
\]
其中右端由 \(U\) 的可测函数演算定义,且等式包含定义域相等.
\end{proposition}

\begin{proof}
若 \(E\subset S^{1}\setminus\{1\}\) 的闭包不含 1,令 \(V_{E}=\operatorname{Range}(\mu_{U}(E))\).在该谱子空间上,\(D(u)\) 和 \(1/(u-1)\) 有界,故比较二次型可知 \(D(U)|_{V_{E}}=D(U|_{V_{E}})\).有界函数演算的乘法性给出
\[
D(U)\psi=i(U+I)(U-I)^{-1}\psi,qquad\psi\in V_{E}.
\]
由定理~\ref{thm:cayley-transform},它与 \(A\) 在 \(V_{E}\) 上一致.把 \(S^{1}\setminus\{1\}\) 分解为一列闭包不含 1 的不交集合,应用例~\ref{ex:direct-sum-essential-selfadjoint} 和两个算子的自伴性,即得全空间上的相等.
\end{proof}

\begin{theorem}[无界自伴算子的投影值谱定理]\label{thm:unbounded-pvm-proof}
定义实轴上的投影值测度
\[
\mu_{A}(E):=\mu_{U}(C(E)),
\label{eq:unbounded-pvm-from-cayley}
\]
其中 \(\mu_{U}\) 是 \(U\) 的投影值测度.则
\[
A=\int_{\mathbb{R}}\lambda\,\mathrm{d}\mu_{A}(\lambda).
\]
\end{theorem}

\begin{proof}
对 \(\psi\) 定义标量测度 \(\mu_{\psi}^{A},\mu_{\psi}^{U}\).由式~\eqref{eq:unbounded-pvm-from-cayley},\(\mu_{\psi}^{A}(E)=\mu_{\psi}^{U}(C(E))\).测度论中的抽象换元定理给出
\[
\int_{\mathbb{R}}\lambda^{2}\,\mathrm{d}\mu_{\psi}^{A}(\lambda)
=\int_{S^{1}\setminus\{1\}}D(u)^{2}\,\mathrm{d}\mu_{\psi}^{U}(u),
\]
因此两个算子的定义域相同;将平方去掉则算子本身也相等.结合命题~\ref{prop:cayley-functional-calculus},得到结论.
\end{proof}

\begin{proof}[定理~\ref{thm:unbounded-spectral-first} 的证明]
存在性由定理~\ref{thm:unbounded-pvm-proof} 得到.若 \(\nu_{A}\) 是另一个满足同样积分公式的投影值测度,关于 \(\nu_{A}\) 计算 Cayley 函数 \(C(A)\),得到同一个酉算子 \(U=(A+iI)(A-iI)^{-1}\).再令 \(\nu_{U}(E)=\nu_{A}(C^{-1}(E))\),则 \(\int u\,\mathrm{d}\nu_{U}(u)=U\).正规谱定理的唯一性给出 \(\nu_{U}=\mu_{U}\),从而 \(\nu_{A}=\mu_{A}\).
\end{proof}

\begin{proof}[定理~\ref{thm:unbounded-spectral-direct-integral} 的证明]
对 \(U=C(A)\) 使用有界正规算子的直积分谱定理.得到关于单位圆的 Hilbert 空间族 \(\mathcal{H}_{u}\) 和测度 \(\mu\),使 \(U\) 对应于 \(s(u)\mapsto us(u)\).点 \(u=1\) 不属于 \(U\) 的特征值;因此 \(\mathcal{H}_{1}\) 或为零,或该单点测度为零,不影响直积分.令 \(\mathcal{K}_{\lambda}:=\mathcal{H}_{C(\lambda)}\),并令 \(\nu(E)=\mu(C(E))\).于是可以把直积分转移到实轴;在这个直积分中,\(A\) 对应于 \(s(\lambda)\mapsto\lambda s(\lambda)\),其定义域正由 \(\int\|\lambda s(\lambda)\|^{2}\,\mathrm{d}\nu<\infty\) 给出.
\end{proof}

无界乘法算子形式的证明与上述证明相同,此处不再重复.

\section{练习}
\label{sec:chapter10-exercises}

\begin{exercise}[一参数群的范数连续性]\label{ex:unitary-group-norm-continuity}
\begin{enumerate}[label=(\alph*)]
\item 若 \(A\) 有界自伴,证明 \(U(t)=e^{itA}\) 在算子范数拓扑下连续;
\item 使用谱定理证明若自伴算子 \(A\) 的谱有界,则 \(A\) 有界;
\item 若 \(A\) 不有界,证明 \(e^{itA}\) 不具有算子范数连续性.
\end{enumerate}
\end{exercise}

\begin{exercise}[动量生成的平移群]\label{ex:momentum-generates-translations}
设 \(P_{j}\) 是第~\ref{sec:basic-unbounded-operators} 节定义的动量算子.证明
\[
(e^{itP_{j}}\psi)(x)=\psi(x+t\hbar e_{j}).
\]
\emph{提示:}先使用 Fourier 变换确定 \(e^{itP_{j}}\psi\) 的 Fourier 变换.
\end{exercise}

\begin{exercise}[强意义下的 Schrödinger 方程]\label{ex:strong-schrodinger-equation}
若 \(A\) 自伴,称 \(\psi(t)\) 强意义下满足
\[
\frac{\mathrm{d}\psi}{\mathrm{d}t}=iA\psi(t)
\]
是指每个 \(\psi(t)\in\operatorname{Dom}(A)\),且
\[
\lim_{h\to0}\left\|\frac{\psi(t+h)-\psi(t)}{h}-iA\psi(t)\right\|=0.
\]
若 \(\psi(t)=e^{itA}\psi_{0}\),证明它强意义下满足方程,当且仅当 \(\psi_{0}\in\operatorname{Dom}(A)\).
\end{exercise}

\begin{exercise}[无界自伴算子的近似特征值]\label{ex:unbounded-almost-eigenvalue}
设 \(A\) 无界自伴,存在非零 \(\psi\in\operatorname{Dom}(A)\) 和 \(\gamma\in\mathbb{R}\),满足 \(\|A\psi-\gamma\psi\|<\varepsilon\|\psi\|\).证明谱中存在 \(\widetilde\gamma\),使 \(|\gamma-\widetilde\gamma|<\varepsilon\).
\end{exercise}

\begin{exercise}[有界算子的 Cayley 变换]\label{ex:bounded-cayley}
若 \(A\) 有界自伴,证明
\[
C(A)=(A+iI)(A-iI)^{-1}
\]
是酉算子,且 1 属于其预解集.证明
\[
A=i(C(A)+I)(C(A)-I)^{-1}.
\]
\end{exercise}

\begin{exercise}[不交换乘积的谱半径]\label{ex:noncommuting-spectral-radius}
给出说明引理~\ref{lem:commuting-spectral-radius-product} 在不假设 \(A,B\) 相交换时不成立的例子.
\end{exercise}

\begin{exercise}[正规矩阵的多项式特征向量]\label{ex:normal-polynomial-eigenvector}
给出正规矩阵 \(A\) 和两个变量多项式 \(p\) 的例子,使 \(p(A,A^{*})\) 的某个特征向量不是 \(A\) 的特征向量.
\begin{note}
尽管如此,定理~\ref{thm:normal-polynomial-spectral-mapping} 的矩阵情形证明表明:若 \(\mu\) 是 \(p(A,A^{*})\) 的特征值,则存在一个同时是 \(A\) 特征向量、且对应 \(\mu\) 的 \(p(A,A^{*})\) 特征向量.
\end{note}
\end{exercise}

\begin{exercise}[不变子空间上的正规限制]\label{ex:normal-restriction}
设 \(W\) 是在 \(A,A^{*}\) 下都不变的闭子空间.
\begin{enumerate}[label=(\alph*)]
\item 证明 \((A|_{W})^{*}=A^{*}|_{W}\);
\item 若 \(A\) 正规,证明 \(A|_{W}\) 正规.
\end{enumerate}
\end{exercise}

\begin{exercise}[有限维不变子空间]\label{ex:finite-dimensional-normal-invariant}
\begin{enumerate}[label=(\alph*)]
\item 设 \(\mathcal{H}\) 有限维,\(A\) 正规,且 \(W\) 在 \(A\) 下不变.证明 \(W\) 在 \(A^{*}\) 下不变;
\item 给出无限维情形中 (a) 不成立的例子.
\end{enumerate}
\end{exercise}

\begin{exercise}[两个近似特征值的比较]\label{ex:compare-almost-eigenvalues}
若同一个非零向量 \(\psi\) 既是 \(A\) 对应 \(\lambda\) 的 \(\varepsilon\)-近似特征向量,又是对应 \(\mu\) 的 \(\delta\)-近似特征向量,证明 \(|\lambda-\mu|<\varepsilon+\delta\).
\end{exercise}
